% ============================================================
% SECTION 3 — COVERAGE GAP OF POOLED EnbPI
% ============================================================
\section{Coverage Gap of Pooled EnbPI Under Seasonal
         Heteroskedasticity}
\label{sec:pooled_gap}

We now characterise the coverage error incurred when EnbPI is applied
with a pooled calibration buffer under
Assumption~\ref{ass:seasonal}.  The main result,
Proposition~\ref{prop:pooled_gap}, shows that the coverage gap of
pooled EnbPI contains a \emph{season-specific bias term} that does
not vanish as $T \to \infty$, in addition to the terms already
identified by \citet{xu2023conformal}.

%------------------------------------------------------------
\subsection{Intra-Stratum Error Structure}
\label{sec:pooled_gap:structure}
%------------------------------------------------------------

Before stating the result, we clarify how the pooled empirical CDF
relates to the seasonal CDF of the prediction target.

Let the training sample be balanced: $T_s = T/\Sper$ for all
$s \in \{1,\ldots,\Sper\}$.\footnote{The balanced assumption
simplifies notation; all results extend to the unbalanced case
by replacing $1/\Sper$ in~\eqref{eq:pooled_cdf} with $T_s/T$.}
The pooled empirical CDF of LOO residuals is
\[
  \hat{F}_T^{\,\text{pool}}(x)
  \;=\;
  \frac{1}{T}\sum_{t=1}^T \Ind{\hat{\epsilon}_t \leq x}.
\]
Under Assumption~\ref{ass:seasonal} and good estimation quality, this
quantity concentrates around its expectation, which is the pooled CDF
$\bar{F}$ defined in~\eqref{eq:pooled_cdf}.  At prediction time
$T+1$ with season $\sstar$, however, the relevant CDF is $F_{\sstar}$,
not $\bar{F}$.

The following lemma establishes the key coupling that drives the gap.

\begin{lemma}[Coupling under seasonal heteroskedasticity]
\label{lem:coupling}
Under Assumption~\ref{ass:seasonal} and the model~\eqref{eq:dgp},
when $X_{T+1} = x_{T+1}$:
\begin{equation}
  \label{eq:coupling}
  \hat{\epsilon}_{T+1}
  \;=\;
  \epsilon_{T+1}
  + \bigl(f(x_{T+1}) - \hat{f}_{-(T+1)}^{\,\phi}(x_{T+1})\bigr),
\end{equation}
where $\epsilon_{T+1} \sim F_{\sstar}$ and the second term on the
right is the pointwise estimation error of the LOO predictor.
\end{lemma}

\begin{proof}
  Equation~\eqref{eq:coupling} is immediate from~\eqref{eq:dgp}:
  $\hat{\epsilon}_{T+1} = y_{T+1} - \hat{f}_{-(T+1)}^{\,\phi}(x_{T+1})
  = f(x_{T+1}) + \epsilon_{T+1} - \hat{f}_{-(T+1)}^{\,\phi}(x_{T+1})$.
\end{proof}

Lemma~\ref{lem:coupling} is identical to equation~(6) of
\citet{xu2023conformal}.  The difference in our setting is that
$\epsilon_{T+1} \sim F_{\sstar}$ (season-specific), while the pooled
buffer $\hat{\boldsymbol\epsilon}$ contains residuals drawn from
all $\Sper$ seasonal distributions.

%------------------------------------------------------------
\subsection{Main Result: Persistent Coverage Gap}
\label{sec:pooled_gap:main}
%------------------------------------------------------------

We now state and prove the main result of this section.

\begin{proposition}[Pooled coverage gap under seasonal
                    heteroskedasticity]
\label{prop:pooled_gap}
Let Assumptions~\ref{ass:xu_a2} and~\ref{ass:seasonal} hold, and
suppose that within each season the errors
$\{\epsilon_t : s(t) = s\}$ are i.i.d.\ according to $F_s$ for
every $s \in \{1,\ldots,\Sper\}$.  Further assume that $F_{\sstar}$
and each $F_s$ are Lipschitz with constant
$L \;=\; \max_{s} L_s = L_G / \min_s \sigma_s$.
Then, for any $T \in \N$, $\alpha \in (0,1)$, and $\beta \in [0,\alpha]$:
%
\begin{align}
  \label{eq:pooled_gap}
  &\bigl|
    \Prob\!\bigl(Y_{T+1} \in \hat{C}_{T+1}^{\,\alpha,\textup{pool}}
               \mid X_{T+1} = x_{T+1}\bigr) - (1-\alpha)
  \bigr| \notag\\
  &\quad\leq\;
  \underbrace{c_0\sqrt{\frac{\log(16T)}{T}}}_{\text{(i) sampling error}}
  \;+\;
  \underbrace{4(L+1)\!\left(\deltaT^{2/3} + \deltaT\right)}_{\text{(ii) estimation error}}
  \;+\;
  \underbrace{2\,\bsstar}_{\text{(iii) seasonal bias}},
\end{align}
%
where $c_0 > 0$ is an absolute constant ($c_0 = 12$ when the errors
are i.i.d., in which case the sharp DKW inequality of
\citealt{massart1990tight} applies; in the seasonal case
Lemma~\ref{lem:dkw_noniid} in Appendix~\ref{app:proofs} supplies the
required non-i.i.d.\ concentration inequality with a slightly larger
absolute constant).
Terms \textup{(i)} and \textup{(ii)} converge to zero as $T\to\infty$
whenever $\deltaT \to 0$.  Term \textup{(iii)} is constant in $T$
and, by Lemma~\ref{lem:bias_properties}(iii), satisfies
$\bsstar > 0$ whenever $G$ is strictly increasing and season
$\sstar$ has the strictly largest or smallest seasonal scale
(in particular, for a genuinely high-volatility season).
\end{proposition}

\begin{proof}
The proof follows the structure of \citet[Theorem~1]{xu2023conformal},
applied under the seasonal model.  We defer the full proof to
Appendix~\ref{app:proofs} and provide here a sketch
highlighting the source of the additional term~(iii).

\medskip\noindent\textit{Proof sketch.}
As in \citet{xu2023conformal}, coverage of the pooled interval
$\hat{C}_{T+1}^{\alpha,\text{pool}}$ is equivalent to
$\betahat \leq \hat{p}_{T+1} \leq 1-\alpha+\betahat$,
where the empirical $p$-value is
$\hat{p}_{T+1} = \hat{F}_T^{\text{pool}}(\hat{\epsilon}_{T+1})$.
Hence:
\begin{align}
  \label{eq:gap_decomp_1}
  &\bigl|\Prob(Y_{T+1} \in \hat{C}_{T+1}^{\alpha,\text{pool}}) - (1-\alpha)\bigr|
  \notag\\
  &= \bigl|\Prob(\betahat \leq \hat{F}_T^{\text{pool}}(\hat{\epsilon}_{T+1})
     \leq 1-\alpha+\betahat)
     - \Prob(\beta \leq F_{\sstar}(\epsilon_{T+1}) \leq 1-\alpha+\beta)\bigr|.
\end{align}
The term $\Prob(\beta \leq F_{\sstar}(\epsilon_{T+1}) \leq 1-\alpha+\beta) = 1-\alpha$
because $F_{\sstar}(\epsilon_{T+1}) \sim \mathrm{Unif}[0,1]$
(Probability Integral Transform, since $\epsilon_{T+1} \sim F_{\sstar}$
which is continuous).

The discrepancy in~\eqref{eq:gap_decomp_1} is then bounded using
the triangle inequality applied to the chain
$\hat{F}_T^{\text{pool}} \to \tilde{F}_T^{\text{pool}} \to \bar{F} \to F_{\sstar}$,
where $\tilde{F}_T^{\text{pool}}(x) = T^{-1}\sum_t \Ind{\epsilon_t \leq x}$
is the pooled empirical CDF of the \emph{true} (unobserved) errors:
\begin{align}
  \label{eq:triangle}
  \bigl|\hat{F}_T^{\text{pool}}(x) - F_{\sstar}(x)\bigr|
  &\leq
  \underbrace{\bigl|\hat{F}_T^{\text{pool}}(x) -
               \tilde{F}_T^{\text{pool}}(x)\bigr|}_{(A)}
  +
  \underbrace{\bigl|\tilde{F}_T^{\text{pool}}(x) -
               \bar{F}(x)\bigr|}_{(B)}
  +
  \underbrace{\bigl|\bar{F}(x) - F_{\sstar}(x)\bigr|}_{(C)}.
\end{align}

\emph{Term (A)} is bounded via Lemma~2 of \citet{xu2023conformal},
using the LOO coupling~\eqref{eq:coupling} and
Assumption~\ref{ass:xu_a2}, yielding a contribution of order
$(L+1)\deltaT^{2/3}$.

\emph{Term (B)}: the pooled empirical CDF $\tilde{F}_T^{\text{pool}}$
is formed from $T$ \emph{independent} (but not identically distributed)
observations, $T_s = T/\Sper$ from each $F_s$, with expectation
$\bar{F}$.  The sharp Dvoretzky--Kiefer--Wolfowitz (DKW) inequality
of \citet{massart1990tight} applies only to i.i.d.\ samples, so it
cannot be invoked directly here.  Instead,
Lemma~\ref{lem:dkw_noniid} in Appendix~\ref{app:proofs} establishes,
via symmetrisation and the bounded-differences inequality
\citep{boucheron2013concentration}, that
\[
  \sup_x\bigl|\tilde{F}_T^{\text{pool}}(x) - \bar{F}(x)\bigr|
  \;\leq\;
  4\sqrt{\frac{\log(16T)}{T}}
\]
with probability at least $1 - (16T)^{-1}$.  The $T$-based rate of
the i.i.d.\ case is therefore preserved; only the absolute constant
changes, which is why term~(i) of~\eqref{eq:pooled_gap} carries the
constant $c_0$ rather than the i.i.d.\ constant $12$.

\emph{Term (C)} equals $\bsstar$ by
Definition~\ref{def:seasonal_bias}, and is independent of $T$.

Assembling the three terms via the argument of
\citet{xu2023conformal} (Lemmas~1--2 and proof of Theorem~1
therein, adapted to the three-step chain~\eqref{eq:triangle},
with the concentration input supplied by
Lemma~\ref{lem:dkw_noniid}) yields~\eqref{eq:pooled_gap}.
The constant in term~(ii) carries over exactly from
\citeauthor{xu2023conformal}'s proof; the absolute constant
$c_0$ in term~(i) results from tracking
Lemma~\ref{lem:dkw_noniid} through that argument.
\end{proof}

\begin{remark}[Comparison with \citealt{xu2023conformal}]
When $\sigma_s = \sigma$ for all $s$ (homoskedasticity), we have
$F_s = F$ for all $s$, so $\bar{F} = F_{\sstar}$ and $\bsstar = 0$;
moreover the errors are then i.i.d., the sharp DKW inequality of
\citet{massart1990tight} applies, and $c_0 = 12$.
Proposition~\ref{prop:pooled_gap} then reduces to
\citet[Theorem~1]{xu2023conformal} exactly, confirming that our
framework nests theirs as a special case.
\end{remark}

\begin{remark}[Relation to Mondrian conformal prediction]
\label{rem:mondrian}
Stratifying conformal intervals by a categorical label is known as
\emph{Mondrian conformal prediction} (CP) \citep{vovk2005mondrian}.
Under exchangeability, Mondrian CP achieves exact conditional
coverage within each stratum.  Proposition~\ref{prop:pooled_gap}
can be viewed as quantifying the loss from \emph{not} applying
Mondrian-type stratification in a time-series setting where
exchangeability fails.  Our Theorem~\ref{thm:strat_coverage}
in Section~\ref{sec:theory} establishes that stratification by
calendar season recovers asymptotic coverage guarantees under
the weaker assumptions appropriate for time-series data.
\end{remark}

%------------------------------------------------------------
\subsection{Discussion}
\label{sec:pooled_gap:discussion}
%------------------------------------------------------------

Proposition~\ref{prop:pooled_gap} identifies a fundamental
limitation of pooled calibration in seasonally heteroskedastic
series that is \emph{not} resolved by increasing the sample size.
Three practical implications follow.

\begin{enumerate}
\item \textbf{Persistent undercoverage in high-volatility seasons.}
  When $\sigma_{\sstar}$ exceeds the remaining seasonal scales
  (a high-volatility season),
  the pooled quantiles underestimate the true dispersion of errors
  in that season, producing intervals that are systematically too
  narrow.  The converse holds for low-volatility seasons, which
  receive wider intervals than necessary.

\item \textbf{Asymptotic invalidity.}
  Even as $T \to \infty$ and terms~(i) and~(ii)
  in~\eqref{eq:pooled_gap} vanish, the \emph{upper bound} remains
  bounded below by $2\bsstar > 0$.  An upper bound alone does not
  prove that the realised gap stays away from zero; however, the
  limiting-coverage computation in Step~6 of
  Appendix~\ref{app:proofs:prop_pooled}
  (equation~\eqref{eq:cov_approx}) shows that pooled coverage
  converges to
  $F_{\sstar}(\bar{F}^{-1}(1-\alpha+\hat\beta))
   - F_{\sstar}(\bar{F}^{-1}(\hat\beta))$,
  which differs from $1-\alpha$ for generic
  $(G, \sigma_1,\ldots,\sigma_\Sper, \alpha)$ whenever
  $\bsstar > 0$, and the simulations of
  Section~\ref{sec:simulations} confirm a persistent per-season
  gap.  In this sense, pooled EnbPI is \emph{asymptotically
  invalid} in the presence of genuine seasonal heteroskedasticity.

\item \textbf{Estimability of the bias.}
  The seasonal bias $\bsstar$ can be estimated consistently from data
  using the season-specific and pooled empirical CDFs of the LOO
  residuals:
  $\hat{b}_{\sstar} = \sup_x |\hat{F}_{T,\sstar}(x) - \hat{F}_T^{\text{pool}}(x)|$,
  where $\hat{F}_{T,\sstar}$ is the empirical CDF restricted to
  season $\sstar$.  This estimator is consistent by the DKW inequality
  applied separately to the two empirical processes.  A practitioner
  can therefore detect seasonal bias before deploying prediction
  intervals.
\end{enumerate}
